% T09 source: Термины и обозначения.tex lines 1--171
\label{sec:terms}

This appendix consolidates the terms and notation used throughout the article.
The explanations provide a concise reference; the formal definitions and
assumptions are introduced in Section~\ref{sec:formal-model}.

\subsection{Core Terms}

\begin{description}[style=nextline]
  \item[Work Unit] An abstract amount of work whose no-AI effort is equal to
  $X$. In the illustrative scenario, work units are linked to estimates in
  story points, although the model does not require a particular task-size
  scale.

  \item[Task] A part of the project that is indivisible at the selected level
  of decomposition. A task has size $Z_i$, operating mode $r$, and normalized
  duration $k_i^{(r)}$. When the level of decomposition changes, one task may
  be replaced by several subtasks.

  \item[Effort] The amount of work expressed in terms of execution time without
  AI\@. In the no-AI baseline with one developer, effort coincides with elapsed
  time; under parallel execution, total effort and the project's elapsed time
  differ.

  \item[Elapsed and Active Time] Elapsed time is measured from the start to the
  completion of a task and includes waiting. Active time counts only intervals
  during which the worker is directly engaged in execution. These metrics, as
  well as task-flow throughput, are not interchangeable.

  \item[Accepted Result] A result that satisfies the specified acceptance
  criterion. The time spent on review, correction, and repeated review is
  included in the task duration; merely completing agent generation does not
  constitute completion of the work.

  \item[No-AI Baseline] Sequential execution of the specified workload by one
  developer without AI\@. Its duration is denoted by $T_h$ and is used as the
  baseline for calculating speedup. The term used in the English-language
  literature is \emph{baseline}.

  \item[Normalized AI-Assisted Duration] The ratio of the AI-assisted task
  duration to the baseline estimate $Z_iX$. It is denoted by $k_i^{(r)}$ and
  depends on the task, tool, and mode. This is a descriptive quantity, not a
  universal property of an AI model or a causal estimate in the absence of a
  separate identification design.

  \item[AI Coding Agent] A software system based on an AI model that can receive
  a task, work with code and tools, and return a result for automated or human
  acceptance. The short form used in the article is \emph{agent}.

  \item[Human--Agent Cell] The organizational unit considered in the article:
  one developer assigns tasks to one or more agents and accepts their results.
  Modeling multiple interacting cells is left for future work.

  \item[Operating Mode] The manner in which a human and an agent interact while
  performing a task. The interactive, delegated, and fully autonomous modes are
  distinguished. The mode characterizes not the product itself but the sequence
  of task specification, execution, review, and correction.

  \item[Interactive Mode] The developer and agent perform a single task jointly,
  sequentially exchanging requests and results. In the limiting case, the human
  component occupies the entire task duration: $h_i=k_i$.

  \item[Delegated Mode] The developer specifies a task, the agent performs it
  autonomously, and the developer then reviews the result and either accepts it
  or returns it for revision.

  \item[Fully Autonomous Mode] A limiting mode in which a prespecified task has
  no mandatory human-attention phase. The model does not represent discovery
  of new tasks or automatic acceptance of results, and this mode is not
  examined in detail. It must not be conflated with an autonomous agent phase
  of a delegated task: such a phase frees the developer during execution, but
  task specification and acceptance of the result still require
  human-attention phases.

  \item[Process Configuration] A complete description of how the work is
  organized: configured parallelism $P$, the assignment of modes to tasks, the
  tool used, and the task-specification and acceptance protocol. When the tool
  and protocol are fixed, the abbreviated notation $\sigma=(P,r)$ is used for a
  common mode, or $\sigma=(P,\rho)$ for different modes. Configurations should
  be compared as a whole because changing any of these elements may change the
  coefficients $k_i^{(r)}$, phase profiles, and overhead.

  \item[Decomposition] The division of a project into tasks and subtasks with
  explicit boundaries and dependencies. Decomposition determines the
  parallelism that can be exploited and the composition of the critical path;
  excessive fragmentation may increase task-specification and acceptance
  overhead.

  \item[Parallelism] The number of AI-assisted execution streams configured to
  be available concurrently. It is denoted by $P$. In scheduling problems, $P$
  is an integer number of processors, whereas in analytical estimates it may be
  interpreted as an effective value averaged over a period. The number of
  agents launched need not coincide with the parallelism actually attainable.

  \item[Dependency Graph] A directed acyclic graph $G=(V,E)$ whose vertices
  correspond to tasks and whose arcs specify precedence relations. The English
  abbreviation is DAG (\emph{directed acyclic graph}).

  \item[Schedule] An assignment of tasks to processors and of their start and
  completion times. A schedule is called feasible if it respects precedence
  relations, phase order, the number of agent streams, local retention of the
  assigned agent slot by a task once started, and the human-resource constraint.

  \item[Task Phase] A continuous stage of task execution with a uniquely
  specified resource type. The model distinguishes autonomous agent phases and
  human-attention phases; the latter require the developer's attention while
  simultaneously retaining the corresponding agent slot.

  \item[Critical Path] The longest path by duration in the dependency graph.
  Its length $L$ provides a lower bound on the project duration: increasing the
  number of agents does not shorten tasks connected sequentially along this
  path.

  \item[Makespan] The time from the start of the project until the completion of
  the final task with an accepted result. This is the principal time metric in
  the article. In scheduling theory, the corresponding term is
  \emph{makespan}.

  \item[Total Work] The sum of the durations of all AI-assisted tasks when they
  are executed in separate streams. It is denoted by $W$. The quantity $W/P$
  provides an idealized estimate under uniform load distribution.

  \item[Agent Component] Intervals of task execution during which the agent
  works autonomously and the developer is free. The normalized duration is
  denoted by $a_i^{(r)}$.

  \item[Human Component] Intervals during which the developer is engaged in
  task specification, reading, review, or directing corrections for the task.
  The normalized duration is denoted by $h_i^{(r)}$.

  \item[Human Resource] The time and attention of a single developer, shared by
  all streams. The resource has unit capacity: human-attention intervals of
  different tasks cannot overlap.

  \item[Coordination Overhead] Additional costs arising from the parallel
  operation of multiple streams. Cross-stream integration overhead is captured
  by the function $C(P)$, while the developer's cost of switching among streams
  is captured by the multiplier $\gamma(P)$.

  \item[Local Blocking of an Agent Slot] The level-M4 rule under which a task,
  once started, retains its assigned agent until its result is accepted. If a
  task requests the attention of a busy developer, the agent stops working and
  remains occupied while waiting; the other agents continue their tasks.
  Blocking one stream does not imply that the entire system stops.

  \item[Queue for Developer Attention] The set of agent requests awaiting
  service by a single developer. Waiting is not included in the productive
  human time $H$, but it increases the occupancy of retained agent slots and may
  cause the actual completion time to exceed the lower bound $B_4$.

  \item[Reuse of a Blocked Stream] Starting another task in place of an agent
  that is waiting for the developer. Neither this arrangement nor the launch of
  an additional priority stream beyond $P$ is considered in the article.

  \item[Speedup] The ratio of the no-AI baseline time to the time of the selected
  AI-assisted configuration. A value greater than one indicates a reduction in
  project duration, a value equal to one indicates no effect on time, and a
  value less than one indicates a slowdown.

  \item[Lower Bound on Time] A value below which the makespan cannot fall. At
  levels M1--M4, the lower bound need not be attainable even under an optimal
  schedule.

  \item[Bottleneck] A task, path, or resource that determines the lower bound on
  time in the configuration under consideration. In the model, the bottleneck
  may be the aggregate load $W/P$, the longest task $M_N$, the critical path
  $L$, or the human resource $\gamma(P)H$.

  \item[Levels M0--M4] Successive refinements of the model. Level M0 assumes
  ideal divisibility of work; M1 accounts for task indivisibility, M2 for the
  dependency graph, M3 redefines durations to account for coordination
  overhead, and M4 introduces phase ordering, a capacity-constrained human
  resource, and local blocking.
\end{description}
